\documentclass[../DoAn.tex]{subfiles}

\begin{document}
\sloppy

Khi bắt đầu đồ án này, tôi nghĩ bài toán quản lý ký túc xá là bài toán đơn
giản --- chủ yếu là số hoá những gì đang làm thủ công. Càng đi sâu, càng nhận
ra mình đã sai. Đằng sau một quyết định phân phòng là hàng chục ràng buộc chính
sách, là kỳ vọng của sinh viên đến từ tỉnh xa chưa từng thấy phòng ở, là nhu
cầu của cán bộ muốn thay đổi chính sách năm sau mà không cần gọi lập trình viên.
Số hoá đúng nghĩa không phải là chuyển Excel lên web --- mà là hiểu đúng từng
bài toán và thiết kế để người dùng thực sự làm chủ được hệ thống. Chương này
nhìn lại toàn bộ hành trình đó, gồm hai phần: Kết luận --- đặt hệ thống trong
tương quan với các giải pháp tương tự, hệ thống hoá kết quả, nhìn nhận hạn chế
và rút ra bài học; và Hướng phát triển --- lộ trình hoàn thiện chiều sâu cùng
những năng lực hoàn toàn mới.

\section{Kết luận}
\label{section:6.1}

Đề tài đã hoàn thành mục tiêu tổng quát đặt ra ở Chương~1: xây dựng một hệ
thống phần mềm tích hợp phục vụ quản lý và phân bổ KTX, kết hợp engine phân
bổ có khả năng cấu hình chính sách linh hoạt với các giao diện người dùng hiện
đại trên cả nền tảng web và di động. Hệ thống không dừng lại ở việc số hoá các
tác vụ thủ công, mà còn bổ sung những năng lực mà quy trình truyền thống không
thể đáp ứng: mô phỏng tác động chính sách trước khi áp dụng, dự báo nhu cầu
chỗ ở theo nhiều kịch bản, và trực quan hoá không gian cư trú bằng công nghệ ba
chiều.

Nhìn lại, có ba quyết định kiến trúc mà tôi tự hào nhất --- không phải vì chúng
phức tạp, mà vì chúng giải quyết đúng vấn đề thực. Quyết định đầu tiên là đưa
sandbox vào lõi nghiệp vụ như một lớp bắt buộc, không phải tính năng tùy chọn:
một khi đã gửi thông báo phân phòng đến hàng trăm sinh viên, không có ``undo''
nào có thể khắc phục được, vì vậy khả năng kiểm chứng trước khi thực thi thật
là yêu cầu thiết kế không thể bỏ qua. Quyết định thứ hai là tách hoàn toàn cấu
hình chính sách khỏi mã nguồn --- không phải để code đẹp hơn, mà là để trao
quyền thay đổi cho đúng người: cán bộ vận hành, không phải lập trình viên.
Quyết định thứ ba là đưa 3D visualization vào hệ thống khi hầu hết mọi người
đều nói ``chỉ cần ảnh tĩnh là đủ'' --- nhưng thực tế đã cho thấy sinh viên từ
tỉnh xa không thể hình dung không gian phòng chỉ qua ảnh, và chính khoảng cách
đó tạo ra làn sóng đổi phòng sau khi nhận.

\subsection{So sánh với các hệ thống tương tự}
\label{subsection:6.1.1}

Để định vị đóng góp của đề tài, hệ thống eDorm được đối chiếu với hai nền tảng
thương mại tiêu biểu trong lĩnh vực chỗ ở sinh viên là Scape \cite{scape_housing}
và University Living \cite{university_living}, cùng với mô hình quản lý KTX
truyền thống đang phổ biến tại nhiều trường đại học Việt Nam dựa trên sổ sách
hoặc cổng thông tin dùng chung. Hai nền tảng quốc tế nói trên có thế mạnh rõ
rệt về trải nghiệm người dùng và luồng đặt phòng trực tuyến tự phục vụ, song
bản chất là các nền tảng thương mại hoá hướng tới thị trường cho thuê, nên không
đặt trọng tâm vào bài toán phân bổ theo chính sách ưu tiên mang tính phúc lợi
như đặc thù của KTX đại học công lập Việt Nam. Bảng~\ref{tab:sosanh} tổng hợp
sự đối chiếu theo chín tiêu chí cốt lõi.

\begin{table}[H]
\centering
\caption{So sánh eDorm với các hệ thống quản lý và đặt phòng ký túc xá tương tự}
\label{tab:sosanh}
\small
\begin{tabular}{p{5.4cm}cccc}
\toprule
\textbf{Tiêu chí} & \textbf{eDorm} & \textbf{Scape} & \textbf{Univ. Living} & \textbf{KTX VN} \\
\midrule
Đăng ký trực tuyến đầu--cuối                 & Có       & Có       & Có       & Không     \\
Phân bổ tự động theo điểm ưu tiên            & Có       & Không    & Không    & Không     \\
Môi trường sandbox (thử $\to$ rollback)      & Có       & Không    & Không    & Không     \\
Dự báo nhu cầu theo 4 kịch bản               & Có       & Không    & Không    & Không     \\
Tham quan ảo 3D phòng (VR tour)              & Có       & Một phần & Một phần & Không     \\
Trực quan hoá địa lý 3D                      & Có       & Không    & Không    & Không     \\
Ứng dụng di động và PWA                      & Có       & Có       & Có       & Không     \\
Thông báo đa kênh (email/SMS/realtime/push)  & Có       & Một phần & Một phần & Một phần  \\
Phù hợp chính sách phân bổ đại học Việt Nam  & Có       & Không    & Không    & Một phần  \\
\bottomrule
\end{tabular}
\end{table}

Kết quả so sánh cho thấy sự vắng mặt của phân bổ theo chính sách ưu tiên trên
cả Scape lẫn University Living không phải là thiếu sót kỹ thuật --- mà là lựa
chọn có chủ đích. Hai nền tảng này là marketplace thương mại: bài toán của họ
là khớp người thuê với chủ nhà theo giá và tiện nghi, không phải xét duyệt
ai xứng đáng được ở dựa trên hoàn cảnh gia đình hay khoảng cách địa lý. Hai bài
toán khác nhau về bản chất, và không thể dùng giải pháp của bài toán này để phán
xét bài toán kia. Ngược lại, nhóm KTX Việt Nam truyền thống lại có logic nghiệp
vụ đúng --- tiêu chí ưu tiên, hạn ngạch theo nhóm năm, ràng buộc giới tính ---
nhưng toàn bộ logic đó đang bị nhốt trong các file Excel và chuỗi email, thiếu
hoàn toàn khả năng kiểm chứng và truy vết. eDorm lấp đầy chính xác khoảng
trống giữa hai nhóm này: giữ nguyên tư duy phúc lợi của KTX đại học Việt Nam,
đồng thời trang bị cho nó nền tảng kỹ thuật hiện đại với đầy đủ khả năng tự
động hoá, mô phỏng và minh bạch hoá quyết định.

\subsection{Những kết quả đã đạt được}
\label{subsection:6.1.2}

Xét trên nhóm nghiệp vụ cốt lõi, hệ thống đã hiện thực hoá trọn vẹn luồng
đăng ký KTX đầu--cuối: từ khâu sinh viên khai báo thông tin ưu tiên và nộp
minh chứng điện tử, qua bước nộp đơn và xét duyệt, đến khâu phân bổ phòng tự
động và phát thông báo kết quả. Trái tim của luồng này là engine phân bổ, vốn
tính điểm ưu tiên trên nhiều tiêu chí --- khoảng cách địa lý, hoàn cảnh gia
đình, diện dân tộc thiểu số (DTTS) và thành tích học tập --- rồi xếp hạng và
gán phòng theo hạn ngạch (quota) của từng nhóm năm học, đồng thời tôn trọng
ràng buộc về giới tính của từng đơn nguyên. Bao quanh engine là môi trường mô
phỏng dạng sandbox: hệ thống nhân bản cơ sở dữ liệu, chạy thử toàn bộ phiên
phân bổ, sinh báo cáo chi tiết cho cán bộ đánh giá, và cho phép khôi phục về
trạng thái ban đầu (rollback) chỉ bằng một thao tác nếu kết quả chưa đạt yêu
cầu. Bên cạnh đó, module dự báo nhu cầu cung cấp bốn kịch bản chỉ tiêu (A, B,
C, D) dựa trên thuật toán lựa chọn linh hoạt giữa tỷ lệ năm gần nhất
(nearest-year, khi chỉ tiêu lệch không quá 5\% so với một năm lịch sử) và
trung bình có trọng số suy giảm theo cấp số nhân ($w = 0{,}7^{\,i}$), kèm
cảnh báo khi tỷ lệ lấp đầy vượt 95\% hoặc xuống dưới 60\%. Hệ thống cũng
quản lý đầy đủ vòng đời cư trú gồm nhận phòng, trả phòng, gia hạn, ghi nhận
vi phạm và thu phí. Sandbox chính là thứ phân biệt eDorm với bất kỳ phần mềm
quản lý KTX nào khác: không chỉ tự động hoá phân bổ, mà còn cho phép cán bộ
kiểm chứng kết quả trước khi áp dụng thật --- điều mà quy trình Excel không bao
giờ làm được mà không đặt dữ liệu sản xuất vào rủi ro.

Xét trên nhóm trải nghiệm người dùng, hệ thống mang tới những hình thức tương
tác vượt ra ngoài khuôn khổ một phần mềm quản lý thông thường. Tính năng tham
quan ảo (Virtual Reality --- VR) ba chiều cho phép sinh viên quan sát phòng ở
ngay trên trình duyệt thông qua thư viện \texttt{Three.js} \cite{threejs} với
mô hình \texttt{glTF/GLB}, hỗ trợ cả chế độ tham quan nhanh theo lộ trình định
sẵn lẫn chế độ khám phá tự do. Ở cấp độ toàn cảnh, năng lực khám phá địa lý
ba chiều dựng trên \texttt{CesiumJS} \cite{cesiumjs} giúp người dùng định vị
các khu KTX trên bản đồ địa cầu với các điểm đánh dấu tương tác. Giao diện
web được xây dựng theo nguyên tắc responsive bằng \texttt{EJS} kết hợp
\texttt{Bootstrap}, trong khi ứng dụng di động phát triển trên
\texttt{React~Native}/\texttt{Expo} \cite{reactnative2024, expo2024router} được
đóng gói qua EAS Build thành tệp cài đặt cho cả Android và iOS. Song song, một
Progressive Web App (PWA) với service worker mang lại khả năng hoạt động ngoại
tuyến và cài đặt trực tiếp từ trình duyệt \cite{webdev_pwa, w3c_manifest}, phục
vụ nhóm người dùng không muốn cài ứng dụng gốc. Hành trình đưa mô hình 3D từ
hơn 200~MB về dưới 8~MB qua ba kỹ thuật tối ưu (Decimate, Texture Baking, Draco)
không chỉ là con số kỹ thuật --- đó là điều kiện tiên quyết để tính năng này
thực sự có thể dùng được trên web với kết nối thực tế của sinh viên Việt Nam.

Xét trên nhóm hạ tầng, hệ thống được đặt trên một nền tảng kỹ thuật chú trọng
cả độ an toàn lẫn khả năng mở rộng. Cơ chế xác thực kết hợp đăng nhập nội bộ
với đăng nhập liên kết qua Google và Microsoft thông qua \texttt{Passport.js}
\cite{passportjs_docs}, sử dụng JSON Web Token (JWT) \cite{rfc7519} cho phiên
di động và bổ sung xác thực hai yếu tố theo chuẩn TOTP (Time-based One-Time
Password) \cite{rfc6238} cho tài khoản quản trị. Hệ thống thông báo vận hành
đa kênh: gửi thư điện tử bằng \texttt{Nodemailer} \cite{nodemailer_docs}, gửi
tin nhắn ngắn qua \texttt{Twilio} \cite{twilio_docs}, và truyền sự kiện thời
gian thực bằng \texttt{Socket.IO} \cite{socketio2024docs} với bộ điều hợp
\texttt{Redis} \cite{redis2024pubsub} cho phép mở rộng theo cụm. Toàn bộ tài
nguyên đa phương tiện --- ảnh phòng, minh chứng và mô hình 360\textdegree{} ---
được lưu trữ trên \texttt{Cloudinary} \cite{cloudinary_docs}. Lớp bảo mật được
gia cố bằng \texttt{Helmet} \cite{helmet_docs}, giới hạn tần suất truy cập
\cite{express_rate_limit}, lọc dữ liệu chống tấn công NoSQL injection, chính
sách bảo mật nội dung (CSP) và băm mật khẩu bằng \texttt{bcrypt}
\cite{bcrypt_npm}. Về vận hành, hệ thống được đóng gói bằng \texttt{Docker}
\cite{docker_docs}, tự động hoá kiểm thử và triển khai qua GitHub Actions
\cite{github_actions}, chạy theo mô hình cụm với \texttt{PM2} \cite{pm2_docs}
sau lớp proxy ngược \texttt{Nginx} \cite{nginx_docs}, và đã được kiểm chứng
vận hành ổn định ở mức khoảng 300 người dùng đồng thời trên cấu hình máy chủ
phổ thông. Con số 300 người dùng đồng thời trên cấu hình phổ thông không phải
là giới hạn trần --- đó là tín hiệu cho thấy kiến trúc đã đủ vững để bàn giao
vận hành thực tế mà không cần đầu tư hạ tầng đặc biệt ngay từ đầu.

\subsection{Hạn chế còn tồn tại}
\label{subsection:6.1.3}

Nhìn nhận hạn chế thẳng thắn quan trọng hơn là che giấu --- một hệ thống tốt
cần biết mình đang đứng ở đâu và còn thiếu gì trước khi đặt mục tiêu tiếp theo.

Module dự báo vẫn còn một khoảng hở đáng kể: phần tính toán số lượng sinh viên
dự kiến rời KTX (expected departures, tương ứng Task~2) chưa được hoàn thiện.
Nguyên nhân không phải kỹ thuật mà là dữ liệu --- hệ thống cần ít nhất ba năm
lịch sử biến động trả phòng để xây dựng mô hình ước lượng đáng tin cậy, và dữ
liệu đó hiện chưa có dưới dạng số hoá tập trung. Kết quả là dự báo hiện chỉ
phản ánh phía nhu cầu đăng ký mới mà chưa cân đối được với dòng sinh viên rời
đi, làm giảm độ chính xác của ước lượng tỷ lệ lấp đầy.

Hệ thống chưa tích hợp cổng thanh toán trực tuyến thực tế như VNPay hay MoMo,
nên các khoản phí KTX vẫn được ghi nhận theo phương thức thủ công thay vì đối
soát tự động. Đây là hạn chế nằm ngoài phạm vi kỹ thuật thuần tuý --- tích hợp
cổng thanh toán đòi hỏi đăng ký pháp lý với đơn vị trung gian thanh toán và
phối hợp với phòng tài chính nhà trường, vượt ra ngoài phạm vi một đồ án tốt
nghiệp.

Hệ thống chưa có lớp giao tiếp dữ liệu (API) đồng bộ với các hệ thống quản lý
đào tạo của nhà trường --- bao gồm hệ thống thông tin sinh viên (SIS) và hệ
thống hoạch định nguồn lực (ERP) --- khiến danh sách sinh viên, điểm học tập và
thông tin học phí vẫn phải nhập liệu hoặc nhập khẩu theo lô. Đây là hạn chế
điển hình của các hệ thống nội bộ đại học: mỗi trường có một cấu trúc SIS riêng,
và việc tích hợp cần đối tác phát triển phía nhà trường cùng đặc tả API mà hiện
chưa có.

Kênh trò chuyện và hỗ trợ trực tuyến giữa sinh viên và ban quản lý chưa được
triển khai, nên các trao đổi phát sinh ngoài quy trình chuẩn --- thắc mắc về
kết quả phân bổ, yêu cầu giải thích cụ thể --- vẫn phải xử lý qua các kênh
ngoài hệ thống. Việc xây dựng kênh này cần cân nhắc kỹ về khối lượng vận hành:
nếu không có cơ chế phân loại và ưu tiên yêu cầu tự động, nguy cơ tạo thêm tải
cho cán bộ thay vì giảm tải.

Cuối cùng, kiểm thử hiện mới dừng ở mức load test thủ công bằng Artillery và
seed data mô phỏng, chưa có user acceptance testing (UAT) với sinh viên và cán
bộ thực tế. Seed data được xây dựng cẩn thận để phản ánh đa dạng trường hợp
biên, nhưng không thể thay thế hoàn toàn cho phản hồi từ người dùng thực khi
họ tương tác với giao diện theo cách mà lập trình viên chưa từng nghĩ đến. UAT
là bước quan trọng không thể bỏ qua trước khi triển khai thật.

\subsection{Bài học kinh nghiệm}
\label{subsection:6.1.4}

Nếu phải chọn một thứ để giữ lại từ toàn bộ quá trình làm đồ án này, tôi sẽ
chọn những bài học thiết kế --- không phải đoạn code nào hay con số nào, mà là
những khoảnh khắc nhận ra mình đã hiểu sai vấn đề và phải làm lại từ đầu.

\textbf{Sandbox là nguyên tắc thiết kế, không phải tính năng.} Lúc đầu, tôi
định để cán bộ chạy thử phân bổ trực tiếp trên dữ liệu thật rồi undo nếu kết
quả chưa ổn. Nghe có vẻ đơn giản --- cho đến khi nhận ra undo trong bối cảnh
phân bổ phòng là không thể. Một khi hệ thống đã gửi email thông báo kết quả
đến 500 sinh viên, không có thao tác kỹ thuật nào có thể ``thu hồi'' được kỳ
vọng đã hình thành. Đó là lúc quyết định sandbox phải là lớp bắt buộc nằm giữa
cán bộ và dữ liệu sản xuất, không phải tuỳ chọn có thể bỏ qua. Bài học ở đây
không phải về kỹ thuật --- mà là về việc hiểu đúng tính không thể đảo ngược
của hành động nghiệp vụ trước khi thiết kế luồng vận hành.

\textbf{Tách cấu hình khỏi code là trao quyền cho đúng người.} Có một câu hỏi
từ phía cán bộ trong quá trình demo đã thay đổi hoàn toàn cách tôi nhìn về thiết
kế hệ thống: ``Năm sau muốn tăng quota năm 1 từ 40\% lên 45\% thì làm thế
nào?''. Câu trả lời lúc đó, nếu logic hard-coded, sẽ là: ``Cần lập trình viên
sửa code, deploy lại, có thể mất vài ngày.'' Đó là câu trả lời không chấp nhận
được trong vận hành thực tế của một tổ chức. Từ câu hỏi đó, tôi đã xây dựng
lại toàn bộ lớp cấu hình để mọi tham số chính sách --- tỷ lệ quota, trọng số
điểm ưu tiên, ngưỡng cảnh báo --- đều có thể thay đổi qua giao diện trong vài
phút, không cần chạm vào mã nguồn. Tách cấu hình khỏi code không phải là quyết
định để code sạch hơn --- đó là quyết định về việc ai thực sự có quyền thay đổi
hành vi của hệ thống.

\textbf{Hiểu nguồn gốc vấn đề trước khi chọn công cụ --- bài học từ 200~MB về
8~MB.} Khi mô hình 3D đầu tiên xuất ra nặng hơn 200~MB, phản xạ đầu tiên là
thử các công cụ nén khác nhau một cách lần lượt. Nhưng cách đó sẽ tốn rất nhiều
thời gian và không chắc ra kết quả. Thay vào đó, tôi dừng lại và phân tích: con
số 200~MB đến từ đâu? Ba nguồn gốc được xác định --- số lượng đa giác quá lớn
trên bề mặt phẳng, quá nhiều texture map riêng lẻ, và dữ liệu hình học chưa
nén. Mỗi nguồn có một công cụ phù hợp: Decimate cho hình học, Texture Baking
cho vật liệu, Draco cho nén. Áp dụng đúng thứ tự, kết quả cuối là dưới 8~MB ---
hơn 25 lần giảm mà không mất chất lượng cảm nhận. Công nghệ đúng mà áp dụng
sai thứ tự thì vẫn không giải quyết được vấn đề; phân tích đúng nguồn gốc mới
là bước đầu tiên, không phải bước cuối cùng.

\textbf{Tầng sự kiện pluggable trả lại giá trị khi scale.} Thiết kế tầng truyền
tải sự kiện có thể hoán đổi giữa cơ chế nội bộ, Redis và Kafka nghe có vẻ over-
engineering lúc đầu --- ứng dụng chỉ chạy trên một instance, cần gì phải phức
tạp vậy? Câu trả lời đến khi thử scale lên hai instance: người dùng ở instance
A không nhận được thông báo do instance B phát ra, vì Socket.IO chưa có bộ đệm
dùng chung. Chuyển sang Redis adapter giải quyết trong một buổi --- nhưng nếu
lúc đầu không thiết kế lớp transport có thể hoán đổi, việc refactor sẽ phức
tạp gấp nhiều lần và có nguy cơ phá vỡ logic nghiệp vụ đang hoạt động. Bài học:
các quyết định kiến trúc về khả năng mở rộng nên được đưa vào từ đầu, dù chưa
cần dùng ngay --- chi phí thiết kế đúng từ đầu luôn thấp hơn chi phí refactor
sau khi hệ thống đã chạy.

\textbf{PWA và React Native bổ trợ nhau, không thay thế nhau.} Ban đầu tôi phân
vân liệu có cần duy trì cả hai hình thức truy cập hay không --- PWA hoặc app
native, chọn một có vẻ đủ và đơn giản hơn. Nhưng quan sát trong quá trình demo
thực tế cho thấy hai nhóm người dùng có hành vi khác nhau rõ rệt: sinh viên có
xu hướng dùng PWA vì không cần cài đặt, truy cập nhanh qua trình duyệt; trong
khi cán bộ thích ứng dụng native vì push notification đáng tin hơn và giao diện
ổn định hơn trên các thao tác phức tạp. Hai nhóm dùng hai hình thức khác nhau
--- duy trì cả hai không phải lãng phí mà là phục vụ đúng nhu cầu của từng nhóm
mà không ép ai phải thay đổi thói quen.

\textbf{Dữ liệu seed đa dạng phát hiện lỗi mà unit test bỏ sót.} Các edge case
của engine phân bổ --- quota đúng bằng số đơn hợp lệ, tie-breaking khi điểm
bằng nhau, gender conflict khi còn đúng một giường, sinh viên có điểm ưu tiên
cực cao nhưng sai giới tính phòng còn trống --- không được phát hiện bởi unit
test mà lộ ra khi chạy seed 500 đơn đa dạng với đủ các nhóm ưu tiên A, B, C,
D. Unit test xác minh từng hàm hoạt động đúng theo spec; seed test đa dạng xác
minh toàn bộ hệ thống phối hợp đúng khi đối mặt với sự phức tạp của thực tế.
Với hệ thống có nhiều ràng buộc tương tác với nhau, hai loại kiểm thử này không
thay thế mà bổ sung cho nhau --- thiếu một trong hai là thiếu một phần bức tranh.

\section{Hướng phát triển}
\label{section:6.2}

Trên nền tảng đã xây dựng, hệ thống có nhiều dư địa để phát triển theo hai
nhánh: hoàn thiện chiều sâu cho những chức năng đã hiện thực và mở rộng chiều
rộng sang những năng lực hoàn toàn mới.

\subsection{Hoàn thiện các chức năng hiện có}
\label{subsection:6.2.1}

Hướng hoàn thiện trước tiên là khép kín module dự báo bằng việc bổ sung phần
ước lượng số lượng sinh viên dự kiến rời KTX khi đã tích luỹ đủ dữ liệu lịch
sử từ ba năm trở lên, qua đó cho phép cân đối đồng thời cả dòng vào và dòng ra
để dự báo tỷ lệ lấp đầy chính xác hơn. Tiếp theo, hệ thống cần tích hợp các
cổng thanh toán phổ biến tại Việt Nam như VNPay, MoMo và ZaloPay để tự động hoá
việc thu và đối soát phí KTX, thay thế hoàn toàn cho thao tác ghi nhận thủ công
hiện nay. Về trải nghiệm trực quan, tính năng tham quan ảo có thể được nâng cấp
lên chuẩn WebXR để hỗ trợ kính thực tế ảo, đồng thời bổ sung các điểm thông tin
(hotspot) về nội thất và chú thích tiện ích ngay trong không gian ba chiều của
phòng. Độ chính xác của dự báo cũng có thể được cải thiện bằng cách đưa thêm
các biến đầu vào có ý nghĩa như tỷ lệ trúng tuyển từng năm và các xu hướng kinh
tế -- xã hội. Một hướng hoàn thiện cụ thể khác là tích hợp xác minh tự động cho
luồng kiểm duyệt minh chứng ưu tiên (priority claims) thông qua API của các đơn
vị cấp giấy tờ liên quan --- hộ nghèo, xác nhận dân tộc thiểu số, giấy tờ học
bổng --- thay vì để cán bộ phải duyệt thủ công từng hồ sơ. Điều này vừa giảm
thời gian xét duyệt, vừa loại bỏ rủi ro giấy tờ giả mạo mà hệ thống hiện tại
chưa kiểm chứng được tự động.

\subsection{Các hướng phát triển mới}
\label{subsection:6.2.2}

Ở chiều mở rộng, hướng đi nền tảng nhất là xây dựng lớp API đồng bộ với hệ
thống SIS và ERP của nhà trường, cho phép tự động cập nhật danh sách sinh viên,
kết quả học tập và tình trạng học phí, qua đó loại bỏ hoàn toàn nhu cầu nhập
liệu thủ công và bảo đảm dữ liệu luôn nhất quán với nguồn gốc. Engine phân bổ
có thể được bổ sung năng lực trí tuệ nhân tạo và học máy (AI/ML), điển hình là
bài toán ghép bạn cùng phòng (roommate matching) dựa trên sở thích và thói quen
sinh hoạt nhằm nâng cao mức độ hài lòng về môi trường ở. Một bảng điều khiển
phân tích kinh doanh thông minh (Business Intelligence --- BI) đa năm sẽ giúp
ban quản lý quan sát bản đồ nhiệt phân bố địa lý của sinh viên, theo dõi xu
hướng cư trú và xuất báo cáo tự động theo từng kỳ. Về hạ tầng, lộ trình mở
rộng theo chiều ngang có thể được hoàn thiện trọn vẹn với điều phối bằng
Kubernetes, phân mảnh dữ liệu trên MongoDB Atlas, cụm \texttt{Redis} và mạng
phân phối nội dung (CDN) cho các mô hình ba chiều dung lượng lớn. Về thiết bị,
hướng KTX thông minh gắn với Internet vạn vật (IoT) hứa hẹn mang lại những
tiện ích thiết thực như nhận phòng tự động bằng mã QR hoặc NFC, theo dõi tiêu
thụ điện nước theo từng phòng và cảnh báo sớm các bất thường trong vận hành.
Cuối cùng, xây dựng module đánh giá mức độ hài lòng sau khi ở (post-stay
feedback) với phân tích sentiment tự động sẽ cho phép dữ liệu phản hồi thực tế
quay lại cải thiện thuật toán ghép phòng và dự báo nhu cầu --- tạo vòng lặp cải
tiến liên tục thay vì chỉ thu thập dữ liệu một chiều rồi để đó.

Khi nhìn lại toàn bộ hành trình từ ý tưởng ban đầu đến hệ thống đang chạy, điều
tôi cảm nhận rõ nhất không phải là những con số kỹ thuật --- không phải 300
người dùng đồng thời, không phải 25 lần nén mô hình 3D, không phải 2 giây phân
bổ. Điều đọng lại là hiểu được rằng phần mềm tốt không đến từ việc chọn đúng
công nghệ, mà đến từ việc hiểu đúng vấn đề của người dùng trước khi mở editor.
Sinh viên đến từ tỉnh xa có thể hình dung được phòng trước khi đăng ký. Cán bộ
quản lý có thể tự điều chỉnh chính sách mà không cần gọi lập trình viên. Mỗi
sinh viên đều biết tại sao mình được hay không được phân phòng. Đó mới là giá
trị thực --- không phải công nghệ phía sau nó.

\end{document}
